\subsection{Physical Interpretation of Signal Masking}
The complete obscuration of the LFM chirp within the distorted waveform $y(t)$ highlights the fundamental challenges inherent in gravitational wave astronomy. The observatory environment acts as a complex convolutional operator that effectively smothers transient astrophysical events. This study models the detection process using the time-domain convolution $y(t)=x(t)*h(t)$, assuming the detector operates as a quasi-stationary Linear Time-Invariant (LTI) system. However, the spectral results confirm significant interference from the 60 Hz mains hum, proving ground-based instruments are physically compromised by the electromagnetic footprint of the electrical grid.

\subsection{Mathematical Mechanics of Noise Enhancement}
The thickened appearance of the raw waveform after spectral inversion in Figure \ref{fig:raw_decon_time} is a manifestation of noise enhancement, an instability inherent in inverse filtering. Division in the frequency domain amplifies noise exponentially where sensitivity $H(f)$ is near zero, characteristic of LIGO detector limits outside the target band. While deconvolution is essential for restoring structural identity, it degrades signal purity, necessitating subsequent filtering.

\subsection{Analytical Comparison of Filter Topologies}
The statistical superiority of the Elliptic filter is interpreted through its aggressive transition band characteristics. Unlike the Butterworth filter, the Elliptic topology utilizes equiripple behavior to achieve a 4th-order transition steepness. This "brick-wall" behavior proved essential for excising 60 Hz power-line harmonics that overlap with the 30-300 Hz band. Sharpness of the transition band proved more critical for signal purity than absolute passband flatness.

\subsection{Phase Distortion and the Failure of Chebyshev II}
The failure of Chebyshev II is attributed to severe non-linear phase distortion. Sharp vertical phase discontinuities observed at 30 Hz and 300 Hz in Figure \ref{fig:bode_plots} indicate a total loss of temporal integrity. Effective reconstruction requires frequency components to be delayed by a consistent amount (linear phase); the Chebyshev II model scrambled the timing of cycles, causing the wave to deconstruct in the time domain.

\subsection{Limitations and Practical Assumptions}
The use of the Bilinear Transformation introduces frequency warping, requiring careful mapping to maintain the 300 Hz cutoff. The assumption of an LTI system is a simplification; in real-world operations, noise is non-stationary and requires adaptive notch filters. Finally, the regularization constant creates a small mathematical bias preventing a perfect 1.0 correlation.


